\documentclass[11pt]{article}
\usepackage[letterpaper,margin=1in]{geometry}
\usepackage[T1]{fontenc}
\usepackage[utf8]{inputenc}
\usepackage{lmodern,microtype}
\usepackage{amsmath,amssymb,mathtools,bm,slashed}
\usepackage{xcolor,enumitem,booktabs,array,needspace}
\usepackage[most]{tcolorbox}
\usepackage{fancyhdr,titlesec,parskip,hyperref}
\definecolor{qftblue}{HTML}{245A8D}
\definecolor{qftgreen}{HTML}{2F6B4F}
\definecolor{qftgray}{HTML}{F4F6F8}
\hypersetup{colorlinks=true,linkcolor=qftblue,urlcolor=qftblue,pdfauthor={Nima Lashkari}}
\pagestyle{fancy}\fancyhf{}
\lhead{PHYS 662: Quantum Field Theory I}\cfoot{\thepage}
\setlength{\headheight}{14pt}
\titleformat{\section}{\Large\bfseries\color{qftblue}}{\thesection.}{0.6em}{}
\titleformat{\subsection}{\large\bfseries\color{qftgreen}}{\thesubsection}{0.6em}{}
\setlist[enumerate]{leftmargin=2em,itemsep=0.25em,topsep=0.35em}
\setlist[itemize]{leftmargin=2em,itemsep=0.2em,topsep=0.35em}
\newcommand{\dd}{\mathrm d}
\newcommand{\ii}{\mathrm i}
\newcommand{\eps}{\varepsilon}
\newcommand{\tr}{\operatorname{tr}}
\newcommand{\Tr}{\operatorname{Tr}}
\newcommand{\diag}{\operatorname{diag}}
\newcommand{\Lag}{\mathcal L}
\newcommand{\D}{\mathcal D}
\newcommand{\vac}{\lvert0\rangle}
\newcommand{\vbra}{\langle0\rvert}
\newcommand{\Tord}{\mathrm T}
\newcommand{\ket}[1]{\lvert#1\rangle}
\newcommand{\bra}[1]{\langle#1\rvert}
\newtcolorbox{keyidea}[1][]{enhanced,breakable,colback=qftgray,colframe=qftblue,
 boxrule=0.7pt,arc=1.5mm,title={#1},fonttitle=\bfseries,coltitle=white}
\newtcolorbox{remarkbox}[1][]{enhanced,breakable,colback=white,colframe=qftgreen,
 boxrule=0.6pt,arc=1.5mm,title={#1},fonttitle=\bfseries,coltitle=white}
\newcommand{\intp}{\int\!\frac{\dd^4p}{(2\pi)^4}}
\newcommand{\MSbar}{\overline{\mathrm{MS}}}
\newcommand{\module}[1]{\clearpage\section{#1}}
\newcommand{\braket}[2]{\langle#1\mid#2\rangle}
\newcommand{\dPi}{\dd\Pi}
\newcommand{\cM}{\mathcal M}
\newcommand{\cH}{\mathcal H}
\newcommand{\cS}{\mathcal S}
\newcommand{\Zp}{Z_{\rm pole}}
\newcommand{\Kallen}{\mathcal K}
\newcommand{\lecturefront}[4]{%
\rhead{Lecture #1}\hypersetup{pdftitle={PHYS 662 Lecture #1: #2}}
\begin{document}\pagenumbering{gobble}
\begin{titlepage}\centering\vspace*{1.05in}
{\scshape\Large PHYS 662 --- Quantum Field Theory I\par}
\vspace{0.35in}{\Huge\bfseries\color{qftblue}Lecture #1\par}
\vspace{0.2in}{\LARGE\bfseries #2\par}
\vspace{0.2in}{\large\color{qftgreen}#3\par}
\vspace{0.6in}{\large Nima Lashkari\par}{\large Fall 2026\par}
\vfill\begin{minipage}{0.86\textwidth}\small #4\par\medskip
Course page: \url{https://nimalashkari.com/quantum-field-theory-2026-part-i/}
\end{minipage}\vspace{0.5in}\end{titlepage}
\pagenumbering{roman}\setcounter{tocdepth}{1}\tableofcontents
\vfill\begin{remarkbox}[Place in the course]
Lecture 26 constructs the physical photon Hilbert space and its propagators.
Lecture 27 develops QED amplitudes. Lecture 28 computes the photon and
fermion self-energies. Lecture 29 relates the vertex to charge
renormalization and the magnetic moment. Homework 9 closes this sequence.
\end{remarkbox}
\small\textbf{Reading and conventions.} The uploaded course slides are the
primary starting point. D.~Tong,
\href{https://www.damtp.cam.ac.uk/user/tong/qft/qfthtml/S6.html}{\textit{Quantum Field Theory}, Chapter 6},
is a companion for photon quantization and QED. For the vertex calculation,
see V.~Kaplunovsky,
\href{https://web2.ph.utexas.edu/~vadim/Classes/2020f/vertex.pdf}{\textit{QED Vertex Correction}},
especially the parameter integral for the Pauli form factor.
For further textbook study, see Peskin--Schroeder, \textit{An Introduction
to QFT}, Chapters 5--7 and 10.
All formulas here use the mostly-minus metric and propagators that include
\(\ii\). Dimensional regularization uses \(d=4-2\eps\), which differs from
some of the slides. The accompanying editorial memo records corrections
and the scope of the supplementary reading.
\normalsize\clearpage\pagenumbering{arabic}}
\lecturefront{28}{Loop Corrections to Propagators}
{Vacuum polarization, screening, and the electron self-energy}
{We compute the two basic QED two-point corrections, keeping their tensor
structure, subtraction scheme, and physical interpretation distinct.
The photon loop leads to running charge and a pair-production threshold;
the electron loop separates mass renormalization from the more delicate
normalization of charged states.}

\section{Which object is transverse?}
The photon propagator \(D_{\mu\nu}\) includes its numerator \(\ii\).
Define the amputated one-particle-irreducible insertion by
\begin{equation}
 \ii\Pi^{\mu\nu}(q)=\ii(q^2\eta^{\mu\nu}-q^\mu q^\nu)\Pi(q^2).
 \label{eq:pidef}
\end{equation}
The Ward identity requires \(q_\mu\Pi^{\mu\nu}=0\). Together with Lorentz
invariance this fixes the tensor structure. It does not require the
full gauge-fixed propagator to be transverse.

For \(q^2\ne0\) let \(P_T=\eta-qq/q^2\) and \(P_L=qq/q^2\).
The free propagator is \(-\ii(P_T+\xi P_L)/q^2\), with Feynman
prescriptions understood. Since the insertion is \(\ii q^2P_T\Pi\),
\begin{align}
 D&=D_0+D_0(\ii\Pi)D_0+D_0(\ii\Pi)D_0(\ii\Pi)D_0+\cdots,\\
 D_{\mu\nu}(q)&=\frac{-\ii P_{T,\mu\nu}}{q^2[1-\Pi_R(q^2)]}
              -\frac{\ii\xi P_{L,\mu\nu}}{q^2}.
 \label{eq:dysonphoton}
\end{align}
Only the transverse part is dressed. The longitudinal sector is fixed
by gauge fixing. Between conserved currents it vanishes, even though
it can contain singular terms at \(q^2=0\).

\begin{keyidea}[A sign convention worth fixing once]
Our tensor is \((q^2\eta^{\mu\nu}-q^\mu q^\nu)\Pi\), and our Dyson
denominator is \(1-\Pi\). A source using the opposite tensor defines
\(\Pi_{\rm alt}=-\Pi\) and has denominator \(1+\Pi_{\rm alt}\).
Change both together. The physical current--current amplitude is unchanged.
\end{keyidea}
For a massive electron loop, \(\Pi_R\) is regular near zero momentum.
The photon remains massless in perturbation theory, while its transverse
residue depends on the normalization convention. On-shell subtraction
sets \(\Pi_R(0)=0\). The local field-renormalization constant \(Z_3\)
should not be confused with the momentum-dependent dressing factor
\([1-\Pi_R(q^2)]^{-1}\).

We use \(\alpha=e^2/(4\pi)\), \(d=4-2\eps\), and
\(1/\bar\eps=1/\eps-\gamma_E+\ln4\pi\). Until the electron self-energy
sections, \(m\) denotes the mass in the one-loop fermion propagator;
using its physical value changes these photon-loop results only beyond
one loop.

\module{A regulator that permits the Ward identity}
The closed electron loop contributes
\begin{equation}
 \ii\Pi^{\mu\nu}_{\rm loop}(q)
 =-e^2\mu^{2\eps}\int\frac{\dd^dk}{(2\pi)^d}
 \frac{\tr[\gamma^\mu(\slashed k+m)\gamma^\nu(\slashed k+\slashed q+m)]}
 {(k^2-m^2+\ii0)[(k+q)^2-m^2+\ii0]}.
 \label{eq:vpint}
\end{equation}
The leading minus sign includes the closed-fermion-loop sign and the
propagator and vertex factors. There is no extra \(1/2\) for this
labeled two-point insertion.

Before evaluating, contract with \(q_\nu\). The identity
\begin{equation}
 \slashed q=(\slashed k+\slashed q-m)-(\slashed k-m)
\end{equation}
converts the integrand into a difference of two single-propagator
traces. Cyclicity and a shift \(k\mapsto k-q\) show that their integrals
cancel. The shift is legitimate only with a regulator that respects
translation of loop momentum. A sharp cutoff applied naively to the
unshifted integral can leave a surface term that looks like a photon
mass. That term is a regulator artifact, not a QED prediction.

Dimensional regularization defines the integral by analytic continuation
from a convergent domain. Our conventions retain \(\tr\mathbf1=4\) and
use \(d\)-dimensional contracted gamma algebra. There is no literal
finite-dimensional spinor representation for arbitrary complex \(d\).
For this vector-current calculation no \(\gamma^5\) prescription is needed.

The scalar integral used repeatedly below is
\begin{equation}
 \mu^{2\eps}\int\frac{\dd^d\ell}{(2\pi)^d}
 \frac1{(\ell^2-\Delta+\ii0)^n}
 =\frac{\ii(-1)^n\mu^{2\eps}}{(4\pi)^{d/2}}
 \frac{\Gamma(n-d/2)}{\Gamma(n)}
 (\Delta-\ii0)^{d/2-n}.
 \label{eq:master}
\end{equation}
It follows by a Wick rotation in a nonsingular domain, a radial Euclidean
integral, and analytic continuation. Continue to physical momenta with
the displayed prescription; do not rotate a contour through poles.

\begin{remarkbox}[What symmetry preservation means here]
Vector QED has no gauge anomaly, and these Ward identities can be
maintained under regulation and subtraction. An anomaly is an obstruction
to preserving the relevant classical symmetry in the quantum theory,
not merely a mistake caused by choosing one inconvenient cutoff.
The axial anomaly is a different question from current conservation here.
\end{remarkbox}

\module{Evaluating the vacuum-polarization tensor}
Combine the denominators using \(1/(AB)=\int_0^1\dd x/[xA+(1-x)B]^2\).
After shifting \(\ell=k+xq\), define
\begin{equation}
 w=x(1-x),\qquad \Delta=m^2-wq^2.
\end{equation}
The gamma trace, with terms odd in \(\ell\) omitted under integration,
is
\begin{equation}
 N^{\mu\nu}=4\left[
 2\ell^\mu\ell^\nu-\eta^{\mu\nu}\ell^2+\eta^{\mu\nu}\Delta
 +2w(q^2\eta^{\mu\nu}-q^\mu q^\nu)\right].
\end{equation}
The apparently nontransverse terms must cancel. To see the cancellation
without guessing, let
\(I_0=\int_\ell(\ell^2-\Delta)^{-2}\) and
\(I_2=\int_\ell\ell^2(\ell^2-\Delta)^{-2}\), with the same regulator.
Lorentz invariance gives
\(\int_\ell\ell^\mu\ell^\nu f(\ell^2)=\eta^{\mu\nu}\int_\ell\ell^2f/d\).
The regulated total derivative
\begin{equation}
 0=\int_\ell\frac{\partial}{\partial\ell_\mu}
 \frac{\ell_\mu}{\ell^2-\Delta}
 =d\int_\ell\frac1{\ell^2-\Delta}-2I_2
\end{equation}
implies \(I_2=d\Delta I_0/(d-2)\). Hence
\begin{equation}
 (2/d-1)I_2+\Delta I_0=0.
\end{equation}
Using \(1/4\) instead of \(1/d\) too early would spoil this step.
The surviving term is manifestly transverse. Applying
\eqref{eq:master} and expanding around four dimensions gives
\begin{equation}
 \boxed{\displaystyle
 \Pi_{\rm loop}(q^2)=-\frac{2\alpha}{\pi}\int_0^1\dd x\,w
 \left[\frac1{\bar\eps}-\ln\frac{m^2-wq^2-\ii0}{\mu^2}\right]}
 +O(\eps).
 \label{eq:piloop}
\end{equation}
In particular \(\Pi_{\rm UV}=-\alpha/(3\pi\bar\eps)\), since
\(\int_0^1w\,\dd x=1/6\). The scalar coefficient has a logarithmic pole;
the quadratic divergence suggested by superficial power counting has
not survived the gauge identity.

The counterterm \(-\delta_3F^2/4\) inserts
\(-\ii\delta_3(q^2\eta^{\mu\nu}-q^\mu q^\nu)\). Therefore
\begin{equation}
 \Pi_R=\Pi_{\rm loop}-\delta_3,\qquad
 \delta_3^{\MSbar}=-\frac{\alpha}{3\pi\bar\eps}.
\end{equation}
These signs are also checked by substituting the insertion back into
the Dyson series, rather than identifying \(\Pi\) from memory.

\module{Subtraction and the effective charge}
The \(\MSbar\) prescription removes only the displayed modified pole.
A physical zero-momentum subtraction instead imposes \(\widehat\Pi(0)=0\):
\begin{equation}
 \widehat\Pi(q^2)=\Pi_{\rm loop}(q^2)-\Pi_{\rm loop}(0)
 =\frac{2\alpha}{\pi}\int_0^1\dd x\,w
 \ln\frac{m^2-wq^2-\ii0}{m^2}.
 \label{eq:pihat}
\end{equation}
For spacelike momentum transfer \(q^2=-Q^2\), this is real and positive.
Exchange between conserved currents is consequently described by
\begin{equation}
 \alpha_{\rm eff}(Q^2)=\frac{\alpha(0)}{1-\widehat\Pi(-Q^2)}.
 \label{eq:alphaeff}
\end{equation}
At fixed one-loop accuracy one may expand the denominator. Retaining it
resums the chain of polarization insertions, but does not include all
higher-order corrections to a process.

Two limits follow directly from the parameter integral. At small
\(|q^2|\), expand the logarithm and use
\(\int_0^1x^2(1-x)^2\dd x=1/30\):
\begin{equation}
 \widehat\Pi(q^2)=-\frac{\alpha}{15\pi}\frac{q^2}{m^2}
 +O(q^4/m^4).
\end{equation}
At \(Q^2\gg m^2\), use
\(\int_0^1w\ln w\,\dd x=-5/18\) to obtain
\begin{equation}
 \widehat\Pi(-Q^2)=\frac{\alpha}{3\pi}
 \left[\ln\frac{Q^2}{m^2}-\frac53\right]
 +O\!\left(\frac{\alpha m^2}{Q^2}\right).
\end{equation}
The effective charge grows as the probe resolves shorter distances.
The low-momentum limit also shows decoupling: a heavy charged particle
contributes only powers of \(Q^2/m^2\) after the long-distance charge is fixed.

\begin{keyidea}[A measured charge and a subtraction parameter]
\(\alpha(0)\) is defined by a long-distance measurement, such as the
Thomson limit. \(\alpha_{\MSbar}(\mu)\) is a parameter in a subtraction
scheme. Their numerical values need not coincide. Matching and running
relate them; physical observables agree when calculated consistently.
\end{keyidea}
At high energies with one active unit-charge Dirac species,
\(\dd\alpha/\dd\ln\mu=2\alpha^2/(3\pi)+O(\alpha^3)\).
Additional species contribute their charge squares when active. A
formal zero of the resummed denominator at exponentially large momentum
is outside controlled weak-coupling perturbation theory; this calculation
does not establish a physical pole there.

\module{The branch cut and pair-production threshold}
For timelike \(q^2=s\), the logarithm in \eqref{eq:pihat} develops an
imaginary part when \(m^2-sx(1-x)<0\). Such an interval exists only for
\(s>4m^2\). Its endpoints are
\begin{equation}
 x_\pm=\frac{1\pm\beta}{2},\qquad
 \beta=\sqrt{1-\frac{4m^2}{s}}.
\end{equation}
Since \(\ln(-a-\ii0)=\ln a-\ii\pi\) for \(a>0\),
\begin{align}
 \operatorname{Im}\widehat\Pi(s+\ii0)
 &=-2\alpha\int_{x_-}^{x_+}\dd x\,x(1-x)\\
 &=-\frac{\alpha}{3}\,
 \beta\left(1+\frac{2m^2}{s}\right)\theta(s-4m^2).
 \label{eq:impi}
\end{align}
The negative sign belongs to our tensor convention. A positive spectral
weight is \(\rho_\Pi(s)=-\operatorname{Im}\widehat\Pi(s+\ii0)/\pi\).
The discontinuity is \(2\ii\operatorname{Im}\widehat\Pi(s+\ii0)\).

The factor \(\beta(1+2m^2/s)\) is the same two-fermion phase-space and
spin-trace factor found in the annihilation cross section. Indeed,
\(\sigma(e^+e^-\to f\bar f)=4\pi\alpha^2\beta(1+2m^2/s)/(3s)\)
when the incoming electron mass is negligible and \(f\) has charge
magnitude \(|e|\). This is the cut of the polarization insertion in
forward scattering, as required by the optical theorem. A real on-shell
photon in vacuum has \(q^2=0\), so it cannot decay into this massive pair.

A once-subtracted dispersion relation reconstructs the analytic function:
\begin{equation}
 \widehat\Pi(q^2)
 =-q^2\int_{4m^2}^\infty\dd s\,
 \frac{\rho_\Pi(s)}{s(s-q^2-\ii0)}.
 \label{eq:dispersion}
\end{equation}
The subtraction implements \(\widehat\Pi(0)=0\). In particular,
\(\widehat\Pi(-Q^2)=Q^2\int\dd s\,\rho_\Pi(s)/[s(s+Q^2)]>0\),
an independent check of screening and of our sign convention.

\begin{remarkbox}[What an imaginary part diagnoses]
A physical cut appears when an allowed intermediate state can go on
shell and pinch the integration contour. It is not enough that one
arbitrarily chosen internal momentum can satisfy its mass-shell equation.
For this two-point function the first such state is a fermion pair;
more complicated amplitudes can have several channel cuts.
\end{remarkbox}

\module{Vacuum polarization as a correction to the potential}
For an electron and a static source of opposite charge magnitude
\(Z|e|\), the leading potential is \(V_C(r)=-Z\alpha/r\).
At zero energy transfer, its momentum-space correction is
\begin{equation}
 \widetilde V(\bm Q)=-\frac{4\pi Z\alpha}{Q^2}
 \left[1+\widehat\Pi(-Q^2)\right]+O(\alpha^3).
\end{equation}
Here \(Z\) labels the source charge, not a field-renormalization constant.
The small-\(Q\) expansion immediately gives
\begin{equation}
 \delta V(\bm r)=-\frac{4Z\alpha^2}{15m^2}\,
 \delta^{(3)}(\bm r)+\hbox{higher derivatives of }\delta^{(3)}(\bm r).
 \label{eq:contact}
\end{equation}
This is a low-momentum effective interaction. It is useful when a bound
wavefunction varies slowly on the scale \(m^{-1}\), and gives
\(\delta E=-4Z\alpha^2|\psi(0)|^2/(15m^2)\) for an S state at this order.
The mass \(m\) is the mass in the vacuum-polarization loop.

To retain the finite range, insert \eqref{eq:dispersion} and use
\begin{equation}
 \int\frac{\dd^3Q}{(2\pi)^3}\,
 \frac{e^{\ii\bm Q\cdot\bm r}}{Q^2+s}
 =\frac{e^{-\sqrt{s}\,r}}{4\pi r}.
\end{equation}
Writing \(s=4m^2t^2\) gives the Uehling potential,
\begin{equation}
 V(r)=-\frac{Z\alpha}{r}\left[1+\frac{2\alpha}{3\pi}
 \int_1^\infty\dd t\,e^{-2mrt}
 \left(1+\frac1{2t^2}\right)\frac{\sqrt{t^2-1}}{t^2}\right].
 \label{eq:uehling}
\end{equation}
The correction increases the magnitude of the attraction relative to
the charge measured at infinite distance. At \(mr\gg1\), the endpoint
\(t=1\) dominates. Expand its weight as
\(3\sqrt{2(t-1)}/2\) and integrate the resulting exponential:
\begin{equation}
 V(r)\sim-\frac{Z\alpha}{r}\left[
 1+\frac{\alpha}{4\sqrt\pi}\frac{e^{-2mr}}{(mr)^{3/2}}
 \left(1+O((mr)^{-1})\right)\right].
\end{equation}
The scale \(2m\) is the pair threshold, not a photon mass. The long-range
Coulomb term survives. The contact expansion and the large-distance tail
are different limits of the same function; a finite number of contact
terms cannot reproduce the tail. Vacuum polarization contributes to
atomic level shifts, but it is only one part of the Lamb shift.

\module{The electron self-energy in Feynman gauge}
Define the electron 1PI insertion as \(-\ii\Sigma(p)\). Its Dyson sum is
\begin{equation}
 S(p)=\frac{\ii}{\slashed p-m-\Sigma_R(p)},\qquad
 \Sigma(p)=A(p^2)\slashed p+B(p^2)m.
 \label{eq:fermiondyson}
\end{equation}
The inverse is a spinor matrix. In Feynman gauge, the loop rules give
\begin{equation}
 -\ii\Sigma_{\rm loop}(p)
 =-e^2\mu^{2\eps}\int\frac{\dd^dk}{(2\pi)^d}
 \frac{\gamma^\mu(\slashed p-\slashed k+m)\gamma_\mu}
 {(k^2+\ii0)[(p-k)^2-m^2+\ii0]}.
\end{equation}
Keep \(p\) off shell initially to separate the ultraviolet calculation
from the on-shell soft-photon singularity. Combining denominators with
parameter \(x\) on the electron denominator and shifting \(k=\ell+xp\)
gives
\begin{equation}
 \Delta_x=xm^2-x(1-x)p^2,\quad
 \gamma^\mu[(1-x)\slashed p-\slashed\ell+m]\gamma_\mu
 =(2-d)[(1-x)\slashed p-\slashed\ell]+dm.
\end{equation}
The odd term integrates to zero. Applying \eqref{eq:master}, with
\(a=\alpha/(4\pi)\), yields
\begin{align}
 \Sigma_{\rm loop}(p)=a\int_0^1\dd x\,\Bigg\{
 &[-2(1-x)\slashed p+4m]
 \left[\frac1{\bar\eps}-\ln\frac{\Delta_x-\ii0}{\mu^2}\right]
 \nonumber\\[-2pt]
 &+2(1-x)\slashed p-2m\Bigg\}.
 \label{eq:selfparameter}
\end{align}
The last two finite terms come from multiplying the \(O(\eps)\) numerator
by the ultraviolet pole. This is a concrete example of why four-dimensional
gamma algebra used too early changes finite answers.

The ultraviolet part is
\begin{equation}
 \Sigma_{\rm UV}(p)=\frac{a}{\bar\eps}(-\slashed p+4m).
\end{equation}
There is no closed-fermion-loop minus sign in this diagram. Its open
fermion line retains the matrix structure \(\slashed p\) and \(m\);
replacing it by a scalar before extracting these coefficients would
lose information needed for both mass and field renormalization.

\module{Mass and field counterterms are different}
Use the renormalized variables of Lecture 27:
\(\psi_0=\sqrt{Z_2}\psi\), \(m_0=Z_m m\), with
\(Z_2=1+\delta_2\), \(Z_m=1+\delta Z_m\).
The quadratic counterterm Lagrangian is
\begin{equation}
 \delta_2\bar\psi(\ii\slashed\partial-m)\psi
 -m\delta Z_m\bar\psi\psi.
\end{equation}
With our insertion convention, the renormalized self-energy is
\begin{equation}
 \Sigma_R=\Sigma_{\rm loop}-\delta_2(\slashed p-m)+m\delta Z_m.
 \label{eq:sigmaren}
\end{equation}
In a general covariant gauge the ultraviolet part of the loop is
\begin{equation}
 \Sigma_{\rm UV}^{(\xi)}
 =\frac{a}{\bar\eps}[-\xi\slashed p+(3+\xi)m].
\end{equation}
Cancellation in \eqref{eq:sigmaren} fixes
\begin{equation}
 \delta_2^{\MSbar}=-\frac{\xi a}{\bar\eps},\qquad
 \delta Z_m^{\MSbar}=-\frac{3a}{\bar\eps}.
 \label{eq:electronct}
\end{equation}
The field constant depends on gauge, while the one-loop mass constant
in this scheme does not. In Landau gauge the field ultraviolet pole
vanishes at this order. That does not make the entire self-energy zero.

Let \(M\) be the perturbative pole mass and \(m(\mu)\) the \(\MSbar\)
mass. To first order, evaluate \(\Sigma_R\) at \(\slashed p=m\) and
\(p^2=m^2\). Equation \eqref{eq:selfparameter} then has
\(\Delta_x=x^2m^2\), so
\begin{align}
 \frac{M-m}{m}
 &=a\int_0^1\dd x\left[
 (2+2x)\left(\ln\frac{\mu^2}{m^2}-2\ln x\right)-2x\right]
 \nonumber\\
 &=a\left[4+3\ln\frac{\mu^2}{m^2}\right].
\end{align}
Consequently,
\begin{equation}
 \boxed{M=m(\mu)\left[1+\frac{\alpha}{4\pi}
 \left(4+3\ln\frac{\mu^2}{m(\mu)^2}\right)\right]+O(\alpha^2).}
\end{equation}
The pole-mass shift is infrared finite at this order and gauge independent.
In an on-shell scheme one instead chooses the finite mass counterterm so
that the parameter equals \(M\). These are two descriptions of the same
mass measurement, not two different electron masses.

\module{Pole normalization and its infrared qualification}
For \(\Sigma_R=A(p^2)\slashed p+B(p^2)m\), the mass condition is
\begin{equation}
 M[1-A(M^2)]-m[1+B(M^2)]=0.
\end{equation}
To extract a regulated pole residue, put \(p\) in its rest frame and
restrict to the positive-energy spinor sector. There the inverse
propagator is \(E[1-A(E^2)]-m[1+B(E^2)]\). Differentiating at \(E=M\)
gives
\begin{equation}
 Z_{\rm pole}^{-1}=1-A(M^2)-2M^2A'(M^2)-2mM B'(M^2).
 \label{eq:residue}
\end{equation}
The derivatives with respect to \(p^2\) are essential. Evaluating only
\(A(M^2)\) would miss part of the residue.

Unlike the mass shift, the on-shell residue is infrared sensitive.
Differentiating the parameter integral probes the endpoint where an
internal photon becomes soft. One must retain an infrared regulator
until external-leg, vertex, and real-emission contributions are combined.
The \(\MSbar\) constant \(Z_2\), which subtracts an ultraviolet pole, is
therefore not automatically the full on-shell residue in \eqref{eq:residue}.
The latter also depends on the gauge and on finite normalization choices.

\begin{keyidea}[Scope of the pole language in QED]
In an infrared-regulated perturbative calculation it is useful to write
\(S(p)\sim\ii Z_{\rm pole}/(\slashed p-M)\). In exact massless-photon QED,
a charged excitation is accompanied by arbitrarily soft radiation and
does not have the isolated one-particle pole assumed in elementary LSZ.
Our pole-mass and residue formulas are perturbative statements with this
infrared qualification; they do not establish an exact charged Fock-state
S-matrix.
\end{keyidea}
For a massive-particle scattering experiment, an inclusive cross section
sums over photons below the detector's resolution. The infrared singular
parts of the virtual correction then cancel against real soft emission.
The finite correction still depends on the stated resolution, as a
measurable probability should.

The two propagators have taught different lessons. Vacuum polarization
renormalizes the transverse photon kinetic term, screens charge, and
records the pair threshold in a branch cut. The electron self-energy
renormalizes both mass and field normalization, with a soft-photon
qualification for the latter. Lecture 29 shows that the vertex is tied
to the electron inverse propagator by the Ward--Takahashi identity, while
its finite Pauli form factor gives a new observable: the anomalous
magnetic moment.
\end{document}
